% ===== PHỤ LỤC A: BẢNG PHÂN CÔNG CÔNG VIỆC =====
\clearpage
\pagestyle{appendixlandscape}
\markboth{}{}


% -------- Trang 1: Tổng quan + Bảng A.1 --------
\begin{landscape}
\chapter{BẢNG PHÂN CÔNG CÔNG VIỆC}

Để bảo đảm tính minh bạch, công bằng và dễ dàng theo dõi tiến độ trong suốt quá trình thực hiện đồ án,
nhóm đã tiến hành phân chia nhiệm vụ một cách rõ ràng ngay từ giai đoạn đầu. Việc phân công cụ thể
giúp xác định trách nhiệm của từng thành viên, giảm thiểu sự chồng chéo trong công việc và đảm bảo
tất cả hoạt động đều được quản lý có hệ thống.

\begin{table}[htbp]
    \caption{Bảng tổng quan phân chia công việc trong nhóm}
    \centering
    \begin{tabularx}{\linewidth}{@{} c c X X c c c @{}}
    \toprule
    \textbf{STT} & \textbf{MSSV} & \textbf{Người thực hiện} & \textbf{Tóm tắt vai trò} &
    \textbf{Số lượng NV} & \textbf{Số lượng hoàn thành đạt} & \textbf{Ghi chú} \\
    \midrule
    1 & 22000001 & Nguyễn Văn A & Trưởng nhóm, Infra & & & \\[4pt]
    2 & 22000002 & Trần Thị B & Lập trình viên, Test & & & \\[4pt]
    3 & 22000003 & Lê Văn C & DevOps, Viết báo cáo & & & \\[4pt]
    4 & 22000004 & Phạm Thị D & Quan sát, Triển khai & & & \\
    \bottomrule
    \end{tabularx}
\end{table}
\end{landscape}

% -------- Trang 2: Bảng A.2 — Nguyễn Văn A --------
\begin{landscape}
\begin{table}[htbp]
    \caption{Phân công công việc chi tiết -- Thành viên Nguyễn Văn A}
    \centering
    \begin{tabularx}{\linewidth}{@{} c X c X c c @{}}
    \toprule
    \textbf{STT} & \textbf{Nội dung nhiệm vụ} & \textbf{Thời gian} & \textbf{Sản phẩm} &
    \textbf{ĐG lần đầu} & \textbf{ĐG lần cuối} \\
    \midrule
    1 & & & & & \\[4pt]
    2 & & & & & \\
    \bottomrule
    \end{tabularx}
\end{table}
\end{landscape}

% -------- Trang 3: Bảng A.3 — Trần Thị B --------
\begin{landscape}
\begin{table}[htbp]
    \caption{Phân công công việc chi tiết -- Thành viên Trần Thị B}
    \centering
    \begin{tabularx}{\linewidth}{@{} c X c X c c @{}}
    \toprule
    \textbf{STT} & \textbf{Nội dung nhiệm vụ} & \textbf{Thời gian} & \textbf{Sản phẩm} &
    \textbf{ĐG lần đầu} & \textbf{ĐG lần cuối} \\
    \midrule
    1 & & & & & \\[4pt]
    2 & & & & & \\
    \bottomrule
    \end{tabularx}
\end{table}
\end{landscape}

% -------- Trang 4: Bảng A.4 + A.5 — Lê Văn C & Phạm Thị D --------
\begin{landscape}
\begin{table}[htbp]
    \caption{Phân công công việc chi tiết -- Thành viên Lê Văn C}
    \centering
    \begin{tabularx}{\linewidth}{@{} c X c X c @{}}
    \toprule
    \textbf{STT} & \textbf{Nội dung nhiệm vụ} & \textbf{Thời gian} & \textbf{Sản phẩm} &
    \textbf{Tự đánh giá} \\
    \midrule
    1 & & & & \\[4pt]
    2 & & & & \\
    \bottomrule
    \end{tabularx}
\end{table}

\vspace{2em}

\begin{table}[htbp]
    \caption{Phân công công việc chi tiết -- Thành viên Phạm Thị D}
    \centering
    \begin{tabularx}{\linewidth}{@{} c X c X c @{}}
    \toprule
    \textbf{STT} & \textbf{Nội dung nhiệm vụ} & \textbf{Thời gian} & \textbf{Sản phẩm} &
    \textbf{Tự đánh giá} \\
    \midrule
    1 & & & & \\[4pt]
    2 & & & & \\
    \bottomrule
    \end{tabularx}
\end{table}
\end{landscape}

\clearpage
\pagestyle{customfancy}   % phục hồi header/footer cho phụ lục B, C
